\PassOptionsToPackage{dvipsnames,svgnames,table}{xcolor}
\documentclass[10pt]{article}
\usepackage[a4paper,margin=22mm]{geometry}
\usepackage[no-math]{fontspec}\setmainfont[SmallCapsFont={Latin Modern Roman Caps}]{Latin Modern Roman}
\usepackage{amsmath,amssymb,amsthm,mathtools,graphicx,xcolor,hyperref,xurl,xspace,ifthen,enumerate,textcomp}
\usepackage[shortlabels]{enumitem}
\setlength{\emergencystretch}{4em}
\newtheorem{theorem}{Theorem}[section]
\newtheorem{theo}[theorem]{Theorem}
\newtheorem{thm}[theorem]{Theorem}
\newtheorem{lemma}[theorem]{Lemma}
\newtheorem{lemm}[theorem]{Lemma}
\newtheorem{prop}[theorem]{Proposition}
\newtheorem{coro}[theorem]{Corollary}
\newtheorem{corollary}[theorem]{Corollary}
\newtheorem{fact}[theorem]{Fact}
\newtheorem{claim}[theorem]{Claim}
\theoremstyle{definition}
\newtheorem{defi}[theorem]{Definition}
\newtheorem{definition}[theorem]{Definition}
\newtheorem{rema}[theorem]{Remark}
\newtheorem{remark}[theorem]{Remark}
\newtheorem{exam}[theorem]{Example}
\newtheorem{example}[theorem]{Example}
\newenvironment{enonce*}[2][]{\par\medskip\noindent\textbf{#2.}\itshape}{\par\medskip}
\providecommand{\eg}{e.g.\ }
\providecommand{\ie}{i.e.\ }
\providecommand{\resp}{resp.\ }
\providecommand{\cf}{cf.\ }
\providecommand{\longthanks}[1]{\par\smallskip\noindent #1\par\smallskip}
\newtheorem{proposition}[theorem]{Proposition}
\newtheorem{conjecture}[theorem]{Conjecture}
\numberwithin{equation}{section}
\newcommand{\R}{{\mathbb R}}
\newcommand{\Z}{{\mathbb Z}}
\newcommand{\N}{{\mathbb N}}
\newcommand{\expref}[2]{\texorpdfstring{\hyperref[#2]{#1~\ref{#2}}}{#1~\ref{#2}}}
\newcommand{\epfmtdoi}[1]{\href{https://doi.org/#1}{doi:#1}}
\newcommand{\epfmt}[2]{#1:\,#2}
\begin{document}
\begingroup\raggedright
\section*{1545. Average-case exponential quantum-query lower bound for cloning hidden-subspace states with primal and dual membership oracles}
Scott Aaronson and Paul Christiano. \emph{Quantum Money from Hidden Subspaces}. Theory of Computing volume 9 article 9 (2013), official complete journal PDF acquired 2026-09-30; canonical DOI 10.4086/toc.2013.v009a009 verified against official landing and printed PDF header. \url{https://theoryofcomputing.org/articles/v009a009/}. CC BY3.0.
\paragraph{Selected problem and scope (editorial).} Exact original Theorem5.5 only: a uniformly random half-dimensional subspace A of F\_2\textasciicircum{}n, one copy of its uniform state, and primal/dual membership oracle access require Omega(sqrt(epsilon) 2\textasciicircum{}\{n/4\}) queries to create a two-register state accepted by the two-copy projector with probability epsilon, for all 1/epsilon=o(2\textasciicircum{}\{n/2\}), with probability over A and the counterfeiter/verifier. Complete new inner-product adversary method, subspace relation/probability calculation, verifier projector, small/high-error reductions, projective counterfeiter amplification and its full hybrid-search argument, source-local fidelity lemma, and invertible-map random self-reduction are included. Standard quantum-state formalism, cited amplitude amplification/fixed-point search and trace-distance/fidelity bound are named prerequisites. Formal serial-number context is conservatively included; cryptographic public-key scheme security, explicit polynomial scheme and Wiesner variants are unselected and uncounted.
\paragraph{Difficulty (editorial).} The target retains the new averaged inner-product progress bound, neighboring-half-subspace relation and exact overlap/probability estimates rather than assuming a no-cloning lower bound as background. Complete new projective amplification and hybrid search convert small-error to high-error cloning, followed by invertible-linear-map random self-reduction for the average-case guarantee. After the named standard quantum-search and fidelity inputs, these source-local query-complexity constructions remain a substantive H2 task beyond routine no-cloning or Grover-search qualifications.
\paragraph{Formatting (editorial).} Exact original human mathematical body and order retained; standalone typography, counter restoration and declared noncore printed locators only. Original comment prose excluded with percent guards retained;1543 author-defined conf=0 journal branch selected with physical newline preservation. No mathematical repair or reconstruction.
\paragraph{Original source quirks (editorial).} Original hybrid-search proof states F(|Phi\_T>,G) at native1169 after introducing Psi\_T; its same printed source form is retained. The random self-reduction paragraph at2517 literally names U\_A in the dual-oracle simulation sentence. All original state/oracle symbols, query bounds and complete surrounding arguments remain unchanged, with no substitution, mathematical refereeing or generated repair.
\endgroup
\clearpage
\setcounter{section}{2}\setcounter{subsection}{1}\setcounter{theorem}{3}
\subsection{Quantum information\label{QI}}

Let us collect a few facts about quantum pure and mixed states that are used
in the paper. We assume basic familiarity with the formalism of bras, kets,
density matrices, etc.; see Nielsen and Chuang~\cite{nc} for a good overview.

Given two mixed states $\rho$ and $\sigma$, their \emph{trace distance} is
defined as 
\[
D\left(  \rho,\sigma\right)  :=\frac{1}{2}\sum_{i=1}^{N}\left\vert
\lambda_{i}\right\vert\,,
\]
where $\lambda_{1},\ldots,\lambda_{N}$ are the
eigenvalues of $\rho-\sigma$. Trace distance is a metric and satisfies
$0\leq D\left(  \rho,\sigma\right)  \leq1$. Also, the \emph{fidelity}
$0\leq F\left(  \rho,\sigma\right)  \leq1$ is defined, in this paper, as the
maximum of $\left\vert \left\langle \psi|\varphi\right\rangle \right\vert
$ over all purifications $\left\vert \psi\right\rangle $ of $\rho$ and
$\left\vert \varphi\right\rangle $ of $\sigma$.\footnote{Some authors instead
define \textquotedblleft fidelity\textquotedblright\ as the maximum of
$\left\vert \left\langle \psi|\varphi\right\rangle \right\vert ^{2}$.} \ By
extension, given a subspace $S$, we let $F\left(  \rho,S\right)  $ be the
maximum of $\left\vert \left\langle \psi|\varphi\right\rangle \right\vert
$ over all purifications $\left\vert \psi\right\rangle $ of $\rho$ and all
unit vectors $\left\vert \varphi\right\rangle \in S$. Trace distance and
fidelity are related as follows~\cite{nc}:

\begin{proposition}
\label{ineq}For all mixed states $\rho,\sigma$,%
\[
D\left(  \rho,\sigma\right)  \leq\sqrt{1-F\left(  \rho,\sigma\right)  ^{2}}\,,
\]
with equality if $\rho$ or $\sigma$ is pure.
\end{proposition}

While fidelity is \emph{not} a metric, it does satisfy the following
inequality, which will be helpful in \expref{Section}{COR}.

\begin{lemma}
[\textquotedblleft Triangle Inequality\textquotedblright\ for Fidelity]%
\label{triangle}Suppose $\left\langle \psi\right\vert \rho\left\vert
\psi\right\rangle \geq1-\varepsilon$ and $\left\langle \varphi\right\vert
\sigma\left\vert \varphi\right\rangle \geq1-\varepsilon$. Then $F\left(
\rho,\sigma\right)  \leq\left\vert \left\langle \psi|\varphi\right\rangle
\right\vert +2\varepsilon^{1/4}$.
\end{lemma}

\begin{proof}
By \expref{Proposition}{ineq},%
\[
D\left(  \rho,\left\vert \psi\right\rangle \right)  \leq\sqrt{1-\left\langle
\psi\right\vert \rho\left\vert \psi\right\rangle }\leq\sqrt{\varepsilon}\,,
\]
and likewise $D\left(  \sigma,\left\vert \varphi\right\rangle \right)
\leq\sqrt{\varepsilon}$. Thus, since trace distance satisfies the triangle
inequality,%
\begin{align*}
D\left(  \rho,\sigma\right)   &  \geq D\left(  \left\vert \psi\right\rangle
,\left\vert \varphi\right\rangle \right)  -D\left(  \rho,\left\vert
\psi\right\rangle \right)  -D\left(  \sigma,\left\vert \varphi\right\rangle
\right)  \\
&  \geq\sqrt{1-\left\vert \left\langle \psi|\varphi\right\rangle \right\vert
^{2}}-2\sqrt{\varepsilon}\,.
\end{align*}
Then%
\begin{align*}
F\left(  \rho,\sigma\right)   &  \leq\sqrt{1-D\left(  \rho,\sigma\right)
^{2}}\\
&  \leq\sqrt{1-\left(  \sqrt{1-\left\vert \left\langle \psi|\varphi
\right\rangle \right\vert ^{2}}-2\sqrt{\varepsilon}\right)  ^{2}}\\
&  \leq\sqrt{\left\vert \left\langle \psi|\varphi\right\rangle \right\vert
^{2}+4\sqrt{\varepsilon}}\\[1ex]
&  \leq\left\vert \left\langle \psi|\varphi\right\rangle \right\vert
+2\varepsilon^{1/4}\,.
\end{align*}

\end{proof}
\setcounter{theorem}{6}
\subsection{Quantum search\label{SEARCH}}

In our security proof for quantum money, an important step will be to
\emph{amplify} a counterfeiter\ who copies a banknote $\$$ with any
non-negligible fidelity to a counterfeiter who copies $\$$ almost perfectly.
\ Taking the contrapositive, this will imply that to rule out the former sort
of counterfeiter, it suffices to rule out the latter.

In this section, we first review two variants of Grover's search algorithm~\cite{grover}\ that are useful for amplifying the fidelity of quantum states.
\ We then introduce a variant that combines the advantages of both.

Assume we are given a pure initial state $\left\vert \operatorname*{Init}%
\right\rangle $, in some Hilbert space $\mathcal{H}$. Our goal is to map
$\left\vert \operatorname*{Init}\right\rangle $ to a final state $\left\vert
\Psi\right\rangle $ that lies in (or close to) a \textquotedblleft good
subspace\textquotedblright\ $G\leq\mathcal{H}$. We have oracle access to two
unitary transformations:

\begin{itemize}
\item $U_{\operatorname*{Init}}$, which maps $\left\vert \operatorname*{Init}%
\right\rangle $ to $-\left\vert \operatorname*{Init}\right\rangle $, and acts
as the identity on all $\left\vert v\right\rangle $ orthogonal to $\left\vert
\operatorname*{Init}\right\rangle $.

\item $U_{G}$, which maps $\left\vert v\right\rangle $ to $-\left\vert
v\right\rangle $ for all $\left\vert v\right\rangle \in G$, and acts as the
identity on all $\left\vert v\right\rangle $ orthogonal to $G$.
\end{itemize}

We are promised that the fidelity of the initial state with $G$,%
\[
F\left(  \left\vert \operatorname*{Init}\right\rangle ,G\right)
=\max_{\left\vert \psi\right\rangle \in G}\left\langle \operatorname*{Init}%
|\psi\right\rangle ,
\]
is at least some $\varepsilon>0$.

In this scenario, \emph{provided }$F\left(  \left\vert \operatorname*{Init}%
\right\rangle ,G\right)  $\emph{\ is known}, the amplitude amplification
framework of Brassard, H\o yer, Mosca, and Tapp~\cite{bhmt} lets us prepare a
state close to $G$ using only $\Theta(1/\varepsilon)$ iterations:

\begin{lemma}
[Amplitude Amplification~\cite{bhmt}]\label{aa}Write $\left\vert
\operatorname*{Init}\right\rangle $ as $\sin\theta\left\vert
\operatorname*{Good}\right\rangle +\cos\theta\left\vert \operatorname*{Bad}%
\right\rangle $, where $\left\vert \operatorname*{Good}\right\rangle $ is the
unit vector formed by projecting $\left\vert \operatorname*{Init}\right\rangle
$ onto $G$, and $\left\vert \operatorname*{Bad}\right\rangle $ is orthogonal
to $\left\vert \operatorname*{Good}\right\rangle $. Then by using $O\left(
T\right)  $ oracle calls\ to $U_{\operatorname*{Init}}$ and $U_{G}$, we can
prepare the state%
\[
\left\vert \Phi_{T}\right\rangle :=\sin\left[  \left(  2T+1\right)
\theta\right]  \left\vert \operatorname*{Good}\right\rangle +\cos\left[
\left(  2T+1\right)  \theta\right]  \left\vert \operatorname*{Bad}%
\right\rangle\,.
\]

\end{lemma}

Note that Grover's algorithm is simply a special case of \expref{Lemma}{aa}, where
$\left\vert \operatorname*{Init}\right\rangle $ is the uniform superposition
over $N$ basis states $\left\vert 1\right\rangle ,\ldots,\left\vert
N\right\rangle $, and $G$ is the subspace spanned by \textquotedblleft
marked\textquotedblright\ states.

However, \expref{Lemma}{aa}\ has an annoying drawback, which it shares with
ordinary Grover search. Namely, the algorithm does \emph{not} converge
monotonically toward the target subspace $G$, but could instead
\textquotedblleft wildly overshoot it,\textquotedblright\ cycling around the
$2$-dimensional\ subspace spanned by $\left\vert \operatorname*{Bad}%
\right\rangle $ and $\left\vert \operatorname*{Good}\right\rangle $. If we
know the fidelity $F(\vert \operatorname*{Init}\rangle
,G)$ in advance (rather than just a lower bound\ on the fidelity),
\emph{or} if we can prepare new copies of $\left\vert \operatorname*{Init}%
\right\rangle $ \textquotedblleft free of charge\textquotedblright\ in case of
failure, then this overshooting is not a serious problem. Alas, neither of
those conditions will hold in our application.

Fortunately, for independent reasons, in 2005 Tulsi, Grover, and Patel~\cite{tgp} introduced a new quantum search algorithm that \emph{does}
guarantee monotonic convergence toward $G$, by alternating unitary
transformations with measurements. (Their algorithm was later simplified and
improved by Chakraborty, Radhakrishnan, and Raghunathan~\cite{crr}.)

\begin{lemma}
[Fixed-Point Quantum Search~\cite{tgp,crr}]\label{fixedpoint}By using $T$
oracle calls\ to $U_{\operatorname*{Init}}$ and $U_{G}$, we can prepare a
state $\left\vert \Psi\right\rangle $ such that $F\left(  \left\vert
\Psi\right\rangle ,G\right)  \geq1-\exp(  -T\varepsilon^{2})$.
\end{lemma}

Rearranging, \expref{Lemma}{fixedpoint}\ lets us prepare a state $\left\vert
\Psi\right\rangle $ such that $F(\vert \Psi\rangle
,G)  \geq1-\delta$ using 
\[
T=O\left(\frac{1}{\varepsilon^{2}}\log
\frac{1}{\delta}\right)
\]
iterations. On the positive side, the dependence
on $1/\delta$ in this bound is logarithmic: we get not only monotonic
convergence toward $G$, but \emph{exponentially-fast} convergence. On the
negative side, notice that the dependence on $\varepsilon$ has worsened from
$1/\varepsilon$ to $1/\varepsilon^{2}$---negating the quadratic speedup that
was the original point of quantum search!

In the rest of this section, we give a \textquotedblleft
hybrid\textquotedblright\ quantum search algorithm that combines the
advantages of Lemmas \ref{aa}\ and \ref{fixedpoint}---i.e., it converges
monotonically toward the target subspace $G$ (rather than \textquotedblleft
overshooting\textquotedblright\ $G$), but \emph{also} achieves a quadratic
speedup. In the context of our security proof for quantum money, this hybrid
algorithm will lead to a quadratically-better (and in fact, tight) lower bound
on the number of queries that a counterfeiter needs to make, compared to what
we would get from using \expref{Lemma}{fixedpoint} alone. While this quadratic
improvement is perhaps only of moderate interest, we include the algorithm in
the hope that it will find other applications.

We first give a technical lemma needed to analyze our algorithm.

\begin{lemma}
\label{interval}For all $L,\beta,\eta,\gamma$, there are at most $\left(
L/\beta+1\right)  \left(  2\eta+1\right)  $ integers $T\in\left\{
0,\ldots,L\right\}$ such that $\left\vert T-\left(  \beta n+\gamma\right)
\right\vert <\eta$ for some integer $n$.
\end{lemma}

\begin{proof}
The real interval $\left[  0,L\right]  $ can intersect at most $L/\beta
+1$ intervals $\left(  \beta n+\gamma-\eta,\beta n+\gamma+\eta\right)  $, and
each such interval can contain at most $2\eta+1$ integer points.
\end{proof}

We now give our hybrid of Lemmas \ref{aa}\ and \ref{fixedpoint}.

\begin{theorem}
[Faster Fixed-Point Search]\label{combined}Let $\delta\geq2\varepsilon$.
\ Then by using \[
O\left(  \frac{\log1/\delta}{\varepsilon\delta^{2}}\right)
\]
oracle calls\ to $U_{\operatorname*{Init}}$ and $U_{G}$, we can prepare a
state $\rho$ such that $F\left(  \rho,G\right)  \geq1-\delta$.
\end{theorem}

\begin{proof}
Let $\xi:=\arcsin\varepsilon$; note that $\varepsilon\leq\xi\leq\frac{\pi}%
{2}\varepsilon$. Also let $L:=\left\lceil 100/\xi\right\rceil $ and
$R:=({25}/{\delta^{2}})\left(  2+\log({1}/{\delta})\right)  $. Then the
algorithm is as follows:

\begin{enumerate}
\item[(1)] Choose an integer $T\in\left\{  0,\ldots,L\right\}  $ uniformly at random.

\item[(2)] Apply $T$ iterations of amplitude amplification with $\left\vert
\operatorname*{Init}\right\rangle $ as the initial state and $G$ as the
target subspace (as in \expref{Lemma}{aa}), to obtain a state $\left\vert \Phi
_{T}\right\rangle $.

\item[(3)] Apply $R$ iterations of fixed-point quantum search with
$\left\vert \Phi_{T}\right\rangle $ as the initial state and $G$ as the
target subspace (as in \expref{Lemma}{fixedpoint}), to obtain a state $\left\vert
\Psi_{T}\right\rangle $.
\end{enumerate}

The final output of the above algorithm is
\[
\rho=\operatorname*{E}_{T\in\left\{  0,\ldots,L\right\}  }\bigl[  \left\vert
\Psi_{T}\right\rangle \left\langle \Psi_{T}\right\vert \bigr]  .
\]
Also, the total number of oracle calls to $U_{\operatorname*{Init}}$ and
$U_{G}$ is%
\[
O\left(  TR\right)  =O\left(  \frac{\log1/\delta}{\varepsilon\delta^{2}%
}\right)  .
\]
(The reason this number scales like $TR$ rather than $T+R$ is that, in step
(3), each time we reflect about the initial state $\left\vert \Phi
_{T}\right\rangle $ we need to rerun step (2). Thus, we need $\Theta\left(
T\right)  $ oracle calls within each of the $R$ iterations.)

By \expref{Lemma}{aa}, after step (2) we have a state $\left\vert \Phi
_{T}\right\rangle $ such that
\[
F\bigl(  \left\vert \Phi_{T}\right\rangle ,G\bigr)  =\bigl\vert \left\langle
\Phi_{T}|\operatorname*{Good}\right\rangle \bigr\vert =\bigl\vert \sin\left[
\left(  2T+1\right)  \xi\right]  \bigr\vert .
\]
So for any $\alpha\in\left(  0,1\right)  $,%
\begin{align*}
\Pr_{T\in\left\{  0,\ldots,L\right\}  }\bigl[  F\left(  \left\vert \Phi
_{T}\right\rangle ,G\right)  <\alpha\bigr]   &  =\Pr_{T\in\left\{
0,\ldots,L\right\}  }\bigl[  \left\vert \sin\left[  \left(  2T+1\right)
\xi\right]  \right\vert <\alpha\bigr]  \\
&  =\Pr_{T\in\left\{  0,\ldots,L\right\}  }\bigl[  \exists n\in\mathbb{Z}%
:\left\vert \left(  2T+1\right)  \xi-\pi n\right\vert <\arcsin\alpha\bigr]
\\
&  \leq\frac{\left(  \frac{L}{\pi/2\xi}+1\right)  \left(  \frac{\arcsin\alpha
}{\xi}+1\right)  }{L+1}\\
&  \leq\frac{2}{\pi}\arcsin\alpha+\frac{2\xi}{\pi}+\frac{\arcsin\alpha}%
{100}+\frac{\xi}{100}\\[1ex]
&  \leq1.02\left(  \alpha+\varepsilon\right)  ,
\end{align*}
where the third line uses \expref{Lemma}{interval}.

Now assume $F(\vert \Phi_{T}\rangle ,G)  \geq\alpha$.
\ Then by \expref{Lemma}{fixedpoint}, after step (3) we have a state $\left\vert
\Psi_{T}\right\rangle $ such that%
\[
F(\vert \Phi_{T}\rangle ,G)  \geq1-\exp\left(
-R\alpha^{2}\right)  .
\]
Let us now make the choice $\alpha:=\delta/5$. Then by the union bound, the
\textquotedblleft average\textquotedblright\ output $\rho=\operatorname*{E}%
_{T}\left[  \left\vert \Psi_{T}\right\rangle \left\langle \Psi_{T}\right\vert
\right]  $ satisfies%
\begin{align*}
1-F\left(  \rho,G\right)   &  \leq1.02\left(  \alpha+\varepsilon\right)
+\exp\left(  -R\alpha^{2}\right) \\
&  \leq0.204\delta+0.51\delta+\exp\left(  -\frac{R\delta^{2}}{25}\right) \\
&  <\delta\,.\qedhere
\end{align*}
\end{proof}
\setcounter{section}{3}\setcounter{subsection}{2}\setcounter{theorem}{1}
\begin{definition}
[Mini-Schemes]\label{minischeme}A (public-key)\textbf{ mini-scheme}
$\mathcal{M}$ consists of two polynomial-time quantum algorithms:

\begin{itemize}
\item $\mathsf{Bank}$, which takes as input a security parameter $0^{n}$, and
probabilistically generates a banknote $\$=\left(  s,\rho_{s}\right)  $, where
$s$ is a classical \textbf{serial number}, and $\rho_{s}$ is a quantum money state.

\item $\mathsf{Ver}$, which takes as input an alleged banknote $%
%
%
\hbox{\rm\rlap/c}%
%
$, and either accepts or rejects.
\end{itemize}

We say $\mathcal{M}$ has \textbf{completeness error} $\varepsilon$ if
$\mathsf{Ver}\left(  \$\right)  $ accepts with probability at least
$1-\varepsilon$ for all valid banknotes $\$$. If $\varepsilon=0$ then
$\mathcal{M}$ has perfect completeness. If, furthermore, $\rho
_{s}=\left\vert \psi_{s}\right\rangle \left\langle \psi_{s}\right\vert $ is
always a pure state, and $\mathsf{Ver}$ simply consists of a projective
measurement onto the rank-$1$ subspace spanned by $\left\vert \psi
_{s}\right\rangle $, then we say $\mathcal{M}$ is \textbf{projective}%
.\footnote{We similarly call a full quantum money scheme projective, if
$\mathsf{Ver}\left(  \$\right)  $ consists of a measurement on one part of
$\$$ in the computational basis, followed by a rank-$1$ projective measurement
on the remaining part.}

Let $\mathsf{Ver}_{2}$ (the \textbf{double verifier}) take as input a single
serial number $s$ as well as two (possibly-entangled) states $\sigma_{1}$ and
$\sigma_{2}$, and accept if and only $\mathsf{Ver}\left(  s,\sigma_{1}\right)
$ and $\mathsf{Ver}\left(  s,\sigma_{2}\right)  $ both accept. We say
$\mathcal{M}$ has \textbf{soundness error} $\delta$ if, given any quantum
circuit $C$ of size $\operatorname*{poly}\left(  n\right)  $ (the
\textbf{counterfeiter}), $\mathsf{Ver}_{2}\left(  s,C\left(  \$\right)
\right)  $ accepts with probability at most $\delta$. Here the probability
is over the banknote $\$$ output by $\mathsf{Bank}\left(  0^{n}\right)  $, as
well as the behavior of $\mathsf{Ver}_{2}$ and $C$.

We call $\mathcal{M}$ \textbf{secure} if it has completeness error $\leq
1/3$ and negligible soundness error.{}
\end{definition}
\setcounter{theorem}{4}
The following result is one of the most useful in the paper. Intuitively, it
says that in projective mini-schemes, a counterfeiter that copies a banknote
with \emph{any} non-negligible fidelity can be \textquotedblleft
amplified\textquotedblright\ to a counterfeiter that copies the banknote
almost \emph{perfectly}---or conversely, that to rule out the former sort of
counterfeiter, it suffices to rule out the latter. The proof makes essential
use of the amplitude amplification results from \expref{Section}{SEARCH}.

\begin{theorem}
[Amplification of Counterfeiters]\label{miniamp}Let $\mathcal{M}=\left(
\mathsf{Bank},\mathsf{Ver}\right)  $ be a projective mini-scheme, and let
$\$=\left(  s,\rho\right)  $ be a valid banknote in $\mathcal{M}$. Suppose
there exists a counterfeiter $C$ that copies $\$$ with probability
$\varepsilon>0$:\ that is,%
\[
\Pr\left[  \mathsf{Ver}_{2}\left(  s,C\left(  \$\right)  \right)  \text{
accepts}\right]  \geq\varepsilon.
\]
Then for all $\delta>0$, there is also a modified counterfeiter $C^{\prime}$
(depending only on $\varepsilon$ and $\delta$, not $\$$), which makes%
\[
O\left(  \frac{\log1/\delta}{\sqrt{\varepsilon}\left(  \sqrt{\varepsilon
}+\delta^{2}\right)  }\right)
\]
queries to $C$, $C^{-1}$, and $\mathsf{Ver}$ and which satisfies%
\[
\Pr\left[  \mathsf{Ver}_{2}\left(  s,C^{\prime}\left(  \$\right)  \right)
\text{ accepts}\right]  \geq1-\delta.
\]

\end{theorem}

\begin{proof}
Write $\$$ as a mixture of pure states:%
\[
\$=\sum p_{i}\left\vert \psi_{i}\right\rangle \left\langle \psi_{i}\right\vert
.
\]
By linearity, clearly it suffices to show that%
\[
\Pr\left[  \mathsf{Ver}_{2}\left(  s,C^{\prime}\left(  \left\vert \psi
_{i}\right\rangle \right)  \right)  \text{ accepts}\right]  \geq1-\delta
\]
for all $i$ such that $p_{i}>0$. We focus on $\left\vert \psi\right\rangle
:=\left\vert \psi_{1}\right\rangle $ without loss of generality.

By assumption, there exists a subspace $S$ such that%
\[
\Pr\left[  \mathsf{Ver}\left(  \rho\right)  \text{ accepts}\right]  =F\left(
\rho,S\right)  ^{2}%
\]
for all $\rho$. Then $F\left(  \$,S\right)  =F\left(  \left\vert
\psi\right\rangle ,S\right)  =1$.

Now, just as $\mathsf{Ver}$ is simply a projector onto $S$, so $\mathsf{Ver}%
_{2}$ is a projector onto $S^{\otimes2}$. Thus%
\[
F\left(  C\left(  \left\vert \psi\right\rangle \right)  ,S^{\otimes2}\right)
\geq\sqrt{\varepsilon}.
\]
So consider performing a fixed-point Grover search, with $C\left(  \left\vert
\psi\right\rangle \right)  $ as the initial state and $S^{\otimes2}$ as the
target subspace. By \expref{Lemma}{fixedpoint}, this will produce a state $\rho
$ such that $F\left(  \rho,S^{\otimes2}\right)  \geq1-\delta$ using
$O\left(({1}/{\varepsilon}) \log ({1}/{\delta})\right)  $ Grover
iterations. Each iteration requires a reflection about $C\left(  \left\vert
\psi\right\rangle \right)  $ and a reflection about $S^{\otimes2}$, which can
be implemented using $O\left(  1\right)  $ queries to $C,C^{-1}$ and
$\mathsf{Ver}$ respectively. Therefore the number of queries to $C,C^{-1}$
and $\mathsf{Ver}$ is $O\left(({1}/{\varepsilon}) \log ({1}/{\delta
})\right)$ as well.

If $\delta$ is large compared to $\varepsilon$, then we can instead use
\expref{Theorem}{combined}, which produces a state $\rho$ such that $F\left(
\rho,S^{\otimes2}\right)  \geq1-\delta$ using 
\[
O\left(  \frac{1}
{\sqrt{\varepsilon}\delta^{2}}\log\frac{1}{\delta}\right)
\]
iterations.
Taking the minimum of the two bounds gives us the claimed bound on query complexity.
\end{proof}
\setcounter{section}{4}\setcounter{subsection}{1}\setcounter{theorem}{0}
\subsection{The method\label{THEMETHOD}}

We now introduce the inner-product adversary method. Let $\mathcal{O}$ be a
set of quantum oracles acting on $n$ qubits each. For each $U\in\mathcal{O}%
$, assume there exists a subspace $S_{U}\leq\mathbb{C}^{2^{n}}$ such that

\begin{enumerate}
\item[(i)] $U\left\vert \psi\right\rangle =-\left\vert \psi\right\rangle
$ for all $\left\vert \psi\right\rangle \in S_{U}$, and

\item[(ii)] $U\left\vert \eta\right\rangle =\left\vert \eta\right\rangle
$ for all $\left\vert \eta\right\rangle \in S_{U}^{\bot}$.
\end{enumerate}

Let $R\subset\mathcal{O}\times\mathcal{O}$ be a symmetric binary relation on
$\mathcal{O}$, with the properties that

\begin{enumerate}
\item[(i)] $\left(  U,U\right)  \notin R$ for all $U\in\mathcal{O}$, and

\item[(ii)] for every $U\in\mathcal{O}$ there exists a $V\in\mathcal{O}%
$ such that $\left(  U,V\right)  \in R$.\
\end{enumerate}

Suppose that for all $U\in\mathcal{O}$ and all $\left\vert \eta\right\rangle
\in S_{U}^{\bot}$, we have%
\[
\operatorname*{E}_{V~:~\left(  U,V\right)  \in R}\left[  F\left(  \left\vert
\eta\right\rangle ,S_{V}\right)  ^{2}\right]  \leq\varepsilon\,,
\]
where $F\left(  \left\vert \eta\right\rangle ,S_{V}\right)  =\max_{\left\vert
\psi\right\rangle \in S_{V}}\left\vert \left\langle \eta|\psi\right\rangle
\right\vert $ is the fidelity between $\left\vert \eta\right\rangle $ and
$S_{V}$. Let $Q$ be a quantum oracle algorithm, and let $Q^{U}$ denote $Q$
run with the oracle $U\in\mathcal{O}$. Suppose $Q^{U}$ begins in the state
$\left\vert \Psi_{0}^{U}\right\rangle $ (possibly already dependent on $U$).
\ Let $\left\vert \Psi_{t}^{U}\right\rangle $ denote the state of $Q^{U}%
$ immediately after the $t^{th}$ query. Also, define a \emph{progress
measure} $p_{t}$ by%
\[
p_{t}:=\operatorname*{E}_{U,V~:~\left(  U,V\right)  \in R}\left[  \left\vert
\left\langle \Psi_{t}^{U}|\Psi_{t}^{V}\right\rangle \right\vert \right]  .
\]
The following lemma bounds how much $p_{t}$ can decrease as the result of a
single query.

\begin{lemma}
[Bound on Progress Rate]\label{innerprod}%
\[
p_{t}\geq p_{t-1}-4\sqrt{\varepsilon}.
\]

\end{lemma}

\begin{proof}
Let $\left\vert \Phi_{t}^{U}\right\rangle $ denote the state of $Q^{U}%
$ immediately \emph{before} the $t^{th}$ query. Then for all $t$, it is
clear that $\left\langle \Phi_{t}^{U}|\Phi_{t}^{V}\right\rangle =\left\langle
\Psi_{t-1}^{U}|\Psi_{t-1}^{V}\right\rangle $: in other words, the unitary
transformations that $Q$ performs in between query steps have no effect on
the inner products. So to prove the lemma, it suffices to show the following
inequality:%
\begin{equation}
\operatorname*{E}_{U,V~:~\left(  U,V\right)  \in R}\left[  \left\vert
\left\langle \Phi_{t}^{U}|\Phi_{t}^{V}\right\rangle \right\vert \right]
-\operatorname*{E}_{U,V~:~\left(  U,V\right)  \in R}\left[  \left\vert
\left\langle \Psi_{t}^{U}|\Psi_{t}^{V}\right\rangle \right\vert \right]
\leq4\sqrt{\varepsilon}\,. \tag{*}%
\end{equation}
Let $\left\{  \left\vert i\right\rangle \right\}  _{i\in\left[  B\right]  }%
$ be an arbitrary orthonormal basis for $Q$'s workspace register. Then we
can write%
\begin{align*}
\left\vert \Phi_{t}^{U}\right\rangle  &  =\sum_{i\in\left[  B\right]  }%
\alpha_{t,i}^{U}\left\vert i\right\rangle \left\vert \Phi_{t,i}^{U}%
\right\rangle \\
&  =\sum_{i\in\left[  B\right]  }\left\vert i\right\rangle \left(  \beta
_{t,i}^{U}\left\vert \eta_{t,i}^{U}\right\rangle +\gamma_{t,i}^{U}\left\vert
\psi_{t,i}^{U}\right\rangle \right)  ,
\end{align*}
where $\left\vert \eta_{t,i}^{U}\right\rangle \in S_{U}^{\bot}$ and
$\left\vert \psi_{t,i}^{U}\right\rangle \in S_{U}$. (By normalization,
$\left\vert \beta_{t,i}^{U}\right\vert ^{2}+\left\vert \gamma_{t,i}%
^{U}\right\vert ^{2}=\left\vert \alpha_{t,i}^{U}\right\vert ^{2}$.) \ A query
transforms the above state to%
\[
\left\vert \Psi_{t}^{U}\right\rangle =\sum_{i\in\left[  B\right]  }\left\vert
i\right\rangle \left(  \beta_{t,i}^{U}\left\vert \eta_{t,i}^{U}\right\rangle
-\gamma_{t,i}^{U}\left\vert \psi_{t,i}^{U}\right\rangle \right)  .
\]
So for all $U,V\in\mathcal{O}$,%
\begin{align*}
\left\langle \Phi_{t}^{U}|\Phi_{t}^{V}\right\rangle -\left\langle \Psi_{t}%
^{U}|\Psi_{t}^{V}\right\rangle  &  =\sum_{i\in\left[  B\right]  }\left(
\overline{\beta}_{t,i}^{U}\left\langle \eta_{t,i}^{U}\right\vert
+\overline{\gamma}_{t,i}^{U}\left\langle \psi_{t,i}^{U}\right\vert \right)
\left(  \beta_{t,i}^{V}\left\vert \eta_{t,i}^{V}\right\rangle +\gamma
_{t,i}^{V}\left\vert \psi_{t,i}^{V}\right\rangle \right) \\
&  ~~~~~~~~~~-\sum_{i\in\left[  B\right]  }\left(  \overline{\beta}_{t,i}%
^{U}\left\langle \eta_{t,i}^{U}\right\vert -\overline{\gamma}_{t,i}%
^{U}\left\langle \psi_{t,i}^{U}\right\vert \right)  \left(  \beta_{t,i}%
^{V}\left\vert \eta_{t,i}^{V}\right\rangle -\gamma_{t,i}^{V}\left\vert
\psi_{t,i}^{V}\right\rangle \right) \\
&  =2\sum_{i\in\left[  B\right]  }\left(  \overline{\beta}_{t,i}^{U}%
\gamma_{t,i}^{V}\left\langle \eta_{t,i}^{U}|\psi_{t,i}^{V}\right\rangle
+\overline{\gamma}_{t,i}^{U}\beta_{t,i}^{V}\left\langle \psi_{t,i}^{U}%
|\eta_{t,i}^{V}\right\rangle \right)  .
\end{align*}
By Cauchy-Schwarz, the above implies that%
\[
\left\vert \left\langle \Phi_{t}^{U}|\Phi_{t}^{V}\right\rangle \right\vert
-\left\vert \left\langle \Psi_{t}^{U}|\Psi_{t}^{V}\right\rangle \right\vert
\leq2\max_{i\in\left[  B\right]  }\left\vert \left\langle \eta_{t,i}^{U}%
|\psi_{t,i}^{V}\right\rangle \right\vert +2\max_{i\in\left[  B\right]
}\left\vert \left\langle \psi_{t,i}^{U}|\eta_{t,i}^{V}\right\rangle
\right\vert .
\]
Now fix $U\in\mathcal{O}$ and $i\in\left[  B\right]  $. Then again applying
Cauchy-Schwarz,%
\begin{align*}
\operatorname*{E}_{V~:~\left(  U,V\right)  \in R}\left[  \left\vert
\left\langle \eta_{t,i}^{U}|\psi_{t,i}^{V}\right\rangle \right\vert \right]
&  \leq\sqrt{\operatorname*{E}_{V~:~\left(  U,V\right)  \in R}\left[
\left\vert \left\langle \eta_{t,i}^{U}|\psi_{t,i}^{V}\right\rangle \right\vert
^{2}\right]  }\\
&  \leq\sqrt{\operatorname*{E}_{V~:~\left(  U,V\right)  \in R}\left[
\max_{\left\vert \psi\right\rangle \in S_{V}}\left\vert \left\langle
\eta_{t,i}^{U}|\psi\right\rangle \right\vert ^{2}\right]  }\\
&  \leq\sqrt{\varepsilon}\,.
\end{align*}
Hence%
\[
\operatorname*{E}_{U,V~:~\left(  U,V\right)  \in R}\left[  \left\vert
\left\langle \eta_{t,i}^{U}|\psi_{t,i}^{V}\right\rangle \right\vert \right]
\leq\sqrt{\varepsilon}%
\]
as well, and likewise%
\[
\operatorname*{E}_{U,V~:~\left(  U,V\right)  \in R}\left[  \left\vert
\left\langle \psi_{t,i}^{U}|\eta_{t,i}^{V}\right\rangle \right\vert \right]
\leq\sqrt{\varepsilon}%
\]
by symmetry. Putting everything together,%
\begin{align*}
p_{t-1}-p_{t}  &  =\operatorname*{E}_{U,V~:~\left(  U,V\right)  \in R}\left[
\left\vert \left\langle \Phi_{t}^{U}|\Phi_{t}^{V}\right\rangle \right\vert
-\left\vert \left\langle \Psi_{t}^{U}|\Psi_{t}^{V}\right\rangle \right\vert
\right] \\
&  \leq2\operatorname*{E}_{U,V~:~\left(  U,V\right)  \in R}\left[  \max
_{i\in\left[  B\right]  }\left\vert \left\langle \eta_{t,i}^{U}|\psi_{t,i}%
^{V}\right\rangle \right\vert \right]  +2\operatorname*{E}_{U,V~:~\left(
U,V\right)  \in R}\left[  \max_{i\in\left[  B\right]  }\left\vert \left\langle
\psi_{t,i}^{U}|\eta_{t,i}^{V}\right\rangle \right\vert \right] \\
&  \leq4\sqrt{\varepsilon}\,.
\end{align*}
This proves inequality (*) and hence the lemma.
\end{proof}

{}From \expref{Lemma}{innerprod}\ we immediately deduce the following.

\begin{theorem}
[Inner-Product Adversary Method]\label{lbthm}Suppose that initially
$\left\vert \left\langle \Psi_{0}^{U}|\Psi_{0}^{V}\right\rangle \right\vert
\geq c$ for all $\left(  U,V\right)  \in R$, whereas by the end we need
$\left\vert \left\langle \Psi_{T}^{U}|\Psi_{T}^{V}\right\rangle \right\vert
\leq d$ for all $\left(  U,V\right)  \in R$. Then $Q$ must make
$T=\Omega(({c-d)}/{\sqrt{\varepsilon}})  $ oracle queries.
\end{theorem}
\setcounter{section}{4}
\section{Classical oracle scheme\label{COR}}

In this section, we construct a mini-scheme, called the \emph{Hidden
Subspace Mini-Scheme}, that requires only a classical oracle. We then use
the inner-product adversary method from Section~4\ to show that our
mini-scheme is secure---indeed, that any counterfeiter must make
$\Omega\left(  2^{n/4}\right)  $ queries to copy a banknote. By the results
of Sections~3.3\ and 2.1, our mini-scheme will
automatically imply a full-blown public-key quantum money scheme, which
requires only a classical oracle and is unconditionally secure.

\subsection{The hidden subspace mini-scheme\label{HSM}}

We identify $n$-bit strings $x\in\left\{  0,1\right\}  ^{n}$ with elements of
the vector space $\mathbb{F}_{2}^{n}$ in the standard way. Then in our
mini-scheme, each $n$-qubit money state will have the form%
\[
\left\vert A\right\rangle :=\frac{1}{\sqrt{\left\vert A\right\vert }}%
\sum_{x\in A}\left\vert x\right\rangle ,
\]
where $A$ is some randomly-chosen subspace of $\mathbb{F}_{2}^{n}$ (i.e., a
set of codewords of a linear code), with $\dim A=n/2$. Let $A^{\bot}$ be
the orthogonal complement of $A$, so that $\dim A^{\bot}=n/2$ as
well. Notice that we can transform $\left\vert A\right\rangle $ to
$\left\vert A^{\bot}\right\rangle $ and vice versa by simply applying
$H_{2}^{\otimes n}$: a Hadamard gate on each of the $n$ qubits, or
equivalently a quantum Fourier transform over $\mathbb{F}_{2}^{n}$.

The basic idea of the mini-scheme is as follows: the bank can easily prepare
the quantum money state $\left\vert A\right\rangle $, starting from a
classical description $\left\langle A\right\rangle $ of $A$ (\eg, a list of
$n/2$ generators). The bank distributes the state $\left\vert
A\right\rangle $, but keeps the classical description $\left\langle
A\right\rangle $ secret. Along with $\left\vert A\right\rangle $ itself,
the bank also publishes details of how to \emph{verify} $\left\vert
A\right\rangle $ by querying two classical oracles, $U_{A}$ and $U_{A^{\bot}%
}$. The first oracle, $U_{A}$, decides membership in $A$: for all $n$-qubit
basis states $\left\vert x\right\rangle $,%
\[
U_{A}\left\vert x\right\rangle =
\begin{cases}
-\left\vert x\right\rangle  & \text{if $x\in A$,}\\
\left\vert x\right\rangle  & \text{otherwise.}%
\end{cases}
\]
The second oracle, $U_{A^{\bot}}$, decides membership in $A^{\bot}$ in the
same way.

Using $U_{A}$, it is easy to implement a projector $\mathbb{P}_{A}$ onto the
set of basis states in $A$. To do so, simply initialize a control qubit to
\[
\left\vert +\right\rangle =\frac{\left\vert 0\right\rangle +\left\vert
1\right\rangle }{\sqrt{2}}\,,
\]
then apply $U_{A}$ conditioned on the control
qubit being in state $\left\vert 1\right\rangle $, then measure the control
qubit in the $\left\{  \left\vert +\right\rangle ,\left\vert -\right\rangle
\right\}  $ basis, and postselect on getting the outcome $\left\vert
-\right\rangle $. Likewise, using $U_{A^{\bot}}$, it is easy to implement a
projector $\mathbb{P}_{A^{\bot}}$ onto the set of basis states in $A^{\bot}$.
\ Then $V_{A}$, the public verification algorithm for the money state
$\left\vert A\right\rangle $, will simply consist of $\mathbb{P}_{A}$, then a
Fourier transform, then $\mathbb{P}_{A^{\bot}}$, and finally a second Fourier
transform to return the legitimate money state back to $\left\vert
A\right\rangle $:%
\[
V_{A}:=H_{2}^{\otimes n}\mathbb{P}_{A^{\bot}}H_{2}^{\otimes n}\mathbb{P}_{A}.
\]
We show in \expref{Lemma}{testworks}\ that $V_{A}$ is just a projector onto
$\left\vert A\right\rangle $. This means, in particular, that $V_{A}%
\left\vert A\right\rangle =\left\vert A\right\rangle $, and that $V_{A}%
$ accepts an arbitrary state $\left\vert \psi\right\rangle $ with
probability $\left\vert \left\langle \psi|A\right\rangle \right\vert ^{2}$.
\ Thus, our mini-scheme is \emph{projective} and has \emph{perfect
completeness}.

But what about security? Intuitively, a counterfeiter could query $U_{A}$ or
$U_{A^{\bot}}$ to find a generating set\ for $A$ or $A^{\bot}$---but that
would require an exponentially-long Grover search, since $\left\vert
A\right\vert =\left\vert A^{\bot}\right\vert =2^{n/2}\ll2^{n}$. Alternatively,
the counterfeiter could measure $\left\vert A\right\rangle $ in the standard
or Hadamard bases---but that would reveal just \emph{one} random element of
$A$ or $A^{\bot}$. Neither ability seems useful for \emph{copying}
$\left\vert A\right\rangle $, let alone recovering a full classical
description of $A$.\footnote{Obviously, if the counterfeiter had
$\Omega(  n)$ copies of $\vert A\rangle $, then it
\emph{could} recover a generating set for $A$, by simply measuring each copy
independently in the standard basis. That is why, in our full quantum money
scheme, the counterfeiter will \emph{not} have $\Omega(n)$ copies of $\vert A\rangle$. Instead, each banknote will
involve a completely different subspace $A_{s}\leq\mathbb{F}_{2}^{n}$
(parameterized by its unique serial number $s$), so that measuring one
banknote reveals nothing about the others.}

And indeed, using the inner-product adversary method plus some other tools, we
will prove the following tight lower bound (\expref{Theorem}{averagecase}): even
if given a single copy of $\left\vert A\right\rangle $, as well as oracle
access to $U_{A}$ and $U_{A^{\bot}}$, a counterfeiter still needs
$\Omega\left(\epsilon2^{n/4}\right)$ queries to prepare a state that
has fidelity $\epsilon$ with $\left\vert A\right\rangle ^{\otimes2}%
$. This will imply that our mini-scheme has $1/\exp(n)$ soundness error.

\subsection{Formal specification\label{SPEC}}

We are not quite done, since we never explained how the bank provides access
to $U_{A}$ and $U_{A^{\bot}}$. Thus, in our \textquotedblleft
final\textquotedblright\ mini-scheme $\mathcal{M}=\left(  \mathsf{Bank}%
_{\mathcal{M}},\mathsf{Ver}_{\mathcal{M}}\right)  $, the bank, verifier, and
counterfeiter will all have access to a \emph{single} classical oracle $U$,
which consists of four components:

\begin{itemize}
\item A \textbf{banknote generator} $\mathcal{G}\left(  r\right)  $, which
takes as input a random string $r\in\left\{  0,1\right\}  ^{n}$, and outputs a
set of linearly independent generators $\langle A_{r}\rangle
=\{  x_{1},\ldots,x_{n/2}\}  $ for a subspace $A_{r}\leq
\mathbb{F}_{2}^{n}$, as well as a unique $3n$-bit \emph{serial number}
$s_{r}\in\left\{  0,1\right\}  ^{3n}$. The function $\mathcal{G}$ is chosen
uniformly at random, subject to the constraint that the serial numbers are all
distinct.\footnote{Note that one can implement $\mathcal{G}$ using an
ordinary random oracle. In that case, the requirement that the serial
numbers are distinct will be satisfied with probability\ $1-O\left(
2^{-n}\right)  $.}

\item A \textbf{serial number checker} $\mathcal{H}\left(  s\right)  $, which
outputs $1$ if $s=s_{r}$ is a valid serial number for some $\left\langle
A_{r}\right\rangle $, and $0$ otherwise.

\item A \textbf{primal subspace tester} $\mathcal{T}_{\operatorname*{primal}}%
$, which takes an input of the form $\left\vert s\right\rangle \left\vert
x\right\rangle $, applies $U_{A_{r}}$ to $\left\vert x\right\rangle $ if
$s=s_{r}$ is a valid serial number for some $\left\langle A_{r}\right\rangle
$, and does nothing otherwise.

\item A \textbf{dual subspace tester} $\mathcal{T}_{\operatorname*{dual}}$,
identical to $\mathcal{T}_{\operatorname*{primal}}$ except that it applies
$U_{A_{r}^{\bot}}$ instead of $U_{A_{r}}$.
\end{itemize}

Then $\mathcal{M}=\left(  \mathsf{Bank}_{\mathcal{M}},\mathsf{Ver}%
_{\mathcal{M}}\right)  $ is defined as follows:

\begin{itemize}
\item $\mathsf{Bank}_{\mathcal{M}}\left(  0^{n}\right)  $ chooses
$r\in\left\{  0,1\right\}  ^{n}$ uniformly at random. It then looks up
$\mathcal{G}\left(  r\right)  =\left(  s_{r},\left\langle A_{r}\right\rangle
\right)  $, and outputs the banknote $\left\vert \$_{r}\right\rangle
=\left\vert s_{r}\right\rangle \left\vert A_{r}\right\rangle $.

\item $\mathsf{Ver}_{\mathcal{M}}\left(
%
%
\hbox{\rm\rlap/c}%
%
\right)  $ first uses $\mathcal{H}$ to check that $%
%
%
\hbox{\rm\rlap/c}%
%
$ has the form $\left(  s,\rho\right)  $, where $s=s_{r}$ is a valid serial
number. If so, then it uses $\mathcal{T}_{\operatorname*{primal}}$ and
$\mathcal{T}_{\operatorname*{dual}}$ to apply $V_{A_{r}}=H_{2}^{\otimes
n}\mathbb{P}_{A_{r}^{\bot}}H_{2}^{\otimes n}\mathbb{P}_{A_{r}}$, and accepts
if and only if $V_{A_{r}}\left(  \rho\right)  $ accepts.
\end{itemize}

\subsection{Analysis\label{ANALYSIS}}

We now analyze the mini-scheme\ defined in Sections \ref{HSM} and 5.2.
\ For convenience, we assume for most of the proof that the subspace
$A\leq\mathbb{F}_{2}^{n}$ is \emph{fixed}, and that the counterfeiter (who
does not know $A$) only has access to the oracles $U_{A}$ and $U_{A^{\bot}}$.
\ Then, at the end, we will explain how to generalize the conclusions to the
\textquotedblleft final\textquotedblright\ mini-scheme $\mathcal{M}$.

It will be convenient to consider the subset $A^{\ast}\subset\left\{
0,1\right\}  ^{n+1}$, defined by%
\[
A^{\ast}:=\left(  0,A\right)  \cup(1,A^{\bot})\,.
\]
Let $S_{A^{\ast}}$ be the subspace of $\mathbb{C}^{2^{n+1}}$ that is spanned
by basis states $\left\vert x\right\rangle $ such that $x\in A^{\ast}$.
\ Then we can think of the pair of oracles $\left(  U_{A},U_{A^{\bot}}\right)
$ as being a \emph{single} oracle $U_{A^{\ast}}$, which satisfies
$U_{A^{\ast}}\left\vert \psi\right\rangle =-\left\vert \psi\right\rangle
$ for all $\left\vert \psi\right\rangle \in S_{A^{\ast}}$, and $U_{A^{\ast}%
}\left\vert \eta\right\rangle =\left\vert \eta\right\rangle $ for all
$\left\vert \eta\right\rangle \in S_{A^{\ast}}^{\bot}$ (where here $\bot
$ means the orthogonal complement in $\mathbb{C}^{2^{n+1}}$, \emph{not} the
orthogonal complement in $\mathbb{F}_{2}^{n}$!).

Recall the definition of the verifier $V_{A}$:%
\[
V_{A}:=H_{2}^{\otimes n}\mathbb{P}_{A^{\bot}}H_{2}^{\otimes n}\mathbb{P}_{A}\,,
\]
where $\mathbb{P}_{A}$ and $\mathbb{P}_{A^{\bot}}$ denote projective
measurements that accept a basis state $\left\vert x\right\rangle $ if and
only if $x$ belongs to $A$ or $A^{\bot}$ respectively. The following lemma
shows that $V_{A}$ \textquotedblleft works,\textquotedblright\ and indeed that
it gives us a projective mini-scheme.

\begin{lemma}
\label{testworks}$V_{A}=\left\vert A\right\rangle \left\langle A\right\vert $
is simply a projector onto $\left\vert A\right\rangle $. So in particular,
\[
\Pr\left[  V_{A}\left(  \left\vert \psi\right\rangle \right)  \text{
accepts}\right]  =\left\vert \left\langle \psi|A\right\rangle \right\vert
^{2}\,.
\]
\end{lemma}

\begin{proof}
It suffices to show that $V_{A}\left\vert A\right\rangle =\left\vert
A\right\rangle $ and that $V_{A}\left\vert \psi\right\rangle =0$ for all
$\left\vert \psi\right\rangle $ orthogonal to $\left\vert A\right\rangle $.
\ First,%
\begin{align*}
V_{A}\left\vert A\right\rangle  &  =H_{2}^{\otimes n}\mathbb{P}_{A^{\bot}%
}H_{2}^{\otimes n}\mathbb{P}_{A}\left\vert A\right\rangle \\
&  =H_{2}^{\otimes n}\mathbb{P}_{A^{\bot}}H_{2}^{\otimes n}\left\vert
A\right\rangle \\
&  =H_{2}^{\otimes n}\mathbb{P}_{A^{\bot}}|A^{\bot}\rangle\\
&  =H_{2}^{\otimes n}|A^{\bot}\rangle\\
&  =\left\vert A\right\rangle .
\end{align*}
Second, if $\left\langle \psi|A\right\rangle =0$ then we can write%
\[
\left\vert \psi\right\rangle =\sum_{x\in2^{n}}{c_{x}\left\vert x\right\rangle
}%
\]
where $\sum_{x\in A}c_{x}=0$. Then%
\begin{align*}
V_{A}\left\vert \psi\right\rangle  &  =H_{2}^{\otimes n}\mathbb{P}_{A^{\bot}%
}H_{2}^{\otimes n}\mathbb{P}_{A}\sum_{x\in2^{n}}{c_{x}\left\vert
x\right\rangle }\\
&  =H_{2}^{\otimes n}\mathbb{P}_{A^{\bot}}H_{2}^{\otimes n}\sum_{x\in A}%
{c_{x}\left\vert x\right\rangle }\\
&  =\frac{1}{\sqrt{2^{n}}}H_{2}^{\otimes n}\mathbb{P}_{A^{\bot}}\sum_{x\in
A}{c_{x}\sum_{y\bot x}\left\vert y\right\rangle }\\
&  =\frac{1}{\sqrt{2^{n}}}H_{2}^{\otimes n}\sum_{y\in A^{\perp}}{\left\vert
y\right\rangle \sum_{x\in A}c_{x}}\\
&  =0\,.
\end{align*}

\end{proof}

We now show that perfect counterfeiting requires exponentially many queries to
$U_{A^{\ast}}$.

\begin{theorem}
[Lower Bound for Perfect Counterfeiting]\label{nocopyc}Given one copy of
$\left\vert A\right\rangle $, as well as oracle access to $U_{A^{\ast}}$, a
counterfeiter needs $\Omega(  2^{n/4})  $ queries to prepare
$\left\vert A\right\rangle ^{\otimes2}$ with certainty (for a worst-case
$\left\vert A\right\rangle $).
\end{theorem}

\begin{proof}
We will apply \expref{Theorem}{lbthm}. Let the set $\mathcal{O}$ contain
$U_{A^{\ast}}$ for every possible subspace $A\leq\mathbb{F}_{2}^{n}$ with
$\dim A=n/2$. Also, put $\left(  U_{A^{\ast}},U_{B^{\ast}}\right)  \in
R$ if and only if $\dim\left(  A\cap B\right)  =n/2-1$. Then given
$U_{A^{\ast}}\in\mathcal{O}$ and $\left\vert \eta\right\rangle \in
S_{A^{\ast}}^{\bot}$, let%
\[
\left\vert \eta\right\rangle =\sum_{x\in\left\{  0,1\right\}  ^{n+1}\setminus
A^{\ast}}\alpha_{x}\left\vert x\right\rangle .
\]
We have%
\begin{align*}
\operatorname*{E}_{U_{B^{\ast}}~:~\left(  U_{A^{\ast}},U_{B^{\ast}}\right)
\in R}\left[  F\left(  \left\vert \eta\right\rangle ,S_{B^{\ast}}\right)
^{2}\right]   &  =\operatorname*{E}_{B~:~\dim\left(  B\right)  =n/2,\dim
\left(  A\cap B\right)  =n/2-1}\left[  \sum_{x\in B^{\ast}\setminus A^{\ast}%
}\left\vert \alpha_{x}\right\vert ^{2}\right] \\
&  \leq\max_{x\in\left\{  0,1\right\}  ^{n+1}\setminus A^{\ast}}\left(
\Pr_{B~:~\dim\left(  B\right)  =n/2,\dim\left(  A\cap B\right)  =n/2-1}\left[
x\in B^{\ast}\right]  \right) \\
&  =\max_{x\in\left\{  0,1\right\}  ^{n}\setminus A}\left(  \Pr_{B~:~\dim
\left(  B\right)  =n/2,\dim\left(  A\cap B\right)  =n/2-1}\left[  x\in
B\right]  \right) \\
&  =\frac{\left\vert B\setminus A\right\vert }{\left\vert \left\{
0,1\right\}  ^{n}\setminus A\right\vert }~~\text{(for }\dim\left(  B\right)
=n/2,~\dim\left(  A\cap B\right)  =n/2-1\text{)}\\
&  =\frac{2^{n/2-1}}{2^{n}-2^{n/2}}\\
&  \leq\frac{1}{2^{n/2}}\,.
\end{align*}
Here the first line uses the definition of fidelity, the second line uses the
easy direction of the minimax theorem, the third line uses the symmetry
between $A$ and $A^{\bot}$, and the fourth line uses the symmetry among all
$2^{n}-2^{n/2}$ strings\ $x\in\left\{  0,1\right\}  ^{n}\setminus A$. The
conclusion is that we can set $\varepsilon:=2^{-n/2}$.

Fix $\left(  U_{A^{\ast}},U_{B^{\ast}}\right)  \in R$. Then $\left\vert
\left\langle A|B\right\rangle \right\vert =1/2$. On the other hand, if the
counterfeiter succeeds, it must map $\left\vert A\right\rangle $ to some
state $\left\vert f_{A}\right\rangle :=\left\vert A\right\rangle \left\vert
A\right\rangle \left\vert \operatorname{garbage}_{A}\right\rangle $, and
$\left\vert B\right\rangle $ to some state $\left\vert f_{B}\right\rangle
:=\left\vert B\right\rangle \left\vert B\right\rangle \left\vert
\operatorname{garbage}_{B}\right\rangle $. Therefore $\left\vert
\left\langle f_{A}|f_{B}\right\rangle \right\vert \leq1/4$. So setting
$c=1/2$ and $d=1/4$, \expref{Theorem}{lbthm}\ tells us that the counterfeiter
must make%
\[
\Omega\left(  \frac{c-d}{\sqrt{\varepsilon}}\right)  =\Omega\left(
2^{n/4}\right)
\]
queries to $U_{A^{\ast}}$.
\end{proof}

A simple modification to the proof of \expref{Theorem}{nocopyc}\ shows that even
to counterfeit money \emph{almost} perfectly, one still needs exponentially
many queries to $U_{A^{\ast}}$.

\begin{corollary}
[Lower Bound for Small-Error Counterfeiting]\label{smallerrror}Given one copy
of $\left\vert A\right\rangle $, as well as oracle access to $U_{A^{\ast}}$, a
counterfeiter needs $\Omega(  2^{n/4})  $ queries to prepare a
state $\rho$ such that $\left\langle A\right\vert ^{\otimes2}\rho\left\vert
A\right\rangle ^{\otimes2}\geq0.9999$ (for a worst-case $\left\vert
A\right\rangle $).
\end{corollary}

\begin{proof}
Let $\left\vert \left\langle A|B\right\rangle \right\vert =c$, and let
$\epsilon=0.0001$. If the counterfeiter succeeds, it must map $\left\vert
A\right\rangle $ to some state $\rho_{A}$, and $\left\vert B\right\rangle
$ to some state $\rho_{B}$, such that $\left\langle A\right\vert ^{\otimes
2}\rho_{A}\left\vert A\right\rangle ^{\otimes2}$ and $\left\langle
B\right\vert ^{\otimes2}\rho_{B}\left\vert B\right\rangle ^{\otimes2}$ are
both at least $1-\epsilon$. So letting $\left\vert f_{A}\right\rangle $ and
$\left\vert f_{B}\right\rangle $ be purifications of $\rho_{A}$ and
$\rho_{B}$ respectively, we have%
\begin{align*}
\left\vert \left\langle f_{A}|f_{B}\right\rangle \right\vert  &  \leq F\left(
\rho_{A},\rho_{B}\right)  \\
&  \leq\left\vert \left\langle A\right\vert ^{\otimes2}\left\vert
B\right\rangle ^{\otimes2}\right\vert +2\epsilon^{1/4}\\
&  =c^{2}+2\epsilon^{1/4}%
\end{align*}
where the second line follows from \expref{Lemma}{triangle}. So setting
$d:=c^{2}+2\epsilon^{1/4}$, \expref{Theorem}{lbthm}\ tells us that the
counterfeiter must make%
\[
\Omega\left(  \frac{c-c^{2}-2\epsilon^{1/4}}{\sqrt{2^{-n/2}}}\right)
\]
queries to $U_{A^{\ast}}$. Fixing $c:=1/2$, the above is $\Omega\left(
2^{n/4}\right)  $.
\end{proof}

Since the verifier $V_{A}$ is projective, we can now combine \expref{Corollary}{smallerrror} with \expref{Theorem}{miniamp} to obtain the following
\textquotedblleft amplified\textquotedblright\ lower bound.

\begin{corollary}
[Lower Bound for High-Error Counterfeiting]\label{largeerrror}Let
$1/\varepsilon=o\left(  2^{n/2}\right)  $. Given one copy of $\left\vert
A\right\rangle $, as well as oracle access to $U_{A^{\ast}}$, a counterfeiter
needs $\Omega\left(  \sqrt{\varepsilon}2^{n/4}\right)  $ queries to prepare a
state $\rho$ such that $\left\langle A\right\vert ^{\otimes2}\rho\left\vert
A\right\rangle ^{\otimes2}\geq\varepsilon$ (for a worst-case $\left\vert
A\right\rangle $).
\end{corollary}

\begin{proof}
Suppose we have a counterfeiter $C$ that makes $o\left(  \sqrt{\varepsilon
}2^{n/4}\right)  $ queries to $U_{A^{\ast}}$, and prepares a state\ $\sigma
$ such that $\left\langle A\right\vert ^{\otimes2}\sigma\left\vert
A\right\rangle ^{\otimes2}\geq\varepsilon$. Let $\delta:=0.00001$. Then by
\expref{Theorem}{miniamp}, there exists an amplified counterfeiter $C^{\prime}%
$ that makes%
\[
O\left(  \frac{\log1/\delta}{\sqrt{\varepsilon}\left(  \sqrt{\varepsilon
}+\delta^{2}\right)  }\right)  =O\left(  \frac{1}{\sqrt{\varepsilon}}\right)
\]
calls to $C$ and $V_{A}$, and that prepares a state $\rho$ such that
$\left\langle A\right\vert ^{\otimes2}\rho\left\vert A\right\rangle
^{\otimes2}\geq1-\delta$. Now, counting the $o\left(  \sqrt{\varepsilon
}2^{n/4}\right)  $ queries from each $C$ invocation and $O\left(  1\right)
$ queries from each $V_{A}$ invocation, the total number of queries that
$C^{\prime}$ makes to $U_{A^{\ast}}$ is%
\[
\left[  o\left(  \sqrt{\varepsilon}2^{n/4}\right)  +O\left(  1\right)
\right]  \cdot O\left(  \frac{1}{\sqrt{\varepsilon}}\right)  =o\left(
2^{n/4}\right)  .
\]
But this contradicts \expref{Corollary}{smallerrror}.
\end{proof}

So far, we have only made statements about the \emph{worst} case for a
would-be counterfeiter. But such guarantees are clearly not enough: it could
be that \emph{most} money states $\left\vert A\right\rangle $ are easy to
duplicate, without contradicting any of the results we have seen so far.

We will show that the problem faced by a counterfeiter is \emph{random
self-reducible}: if a counterfeiter could duplicate a uniformly-random money
state $\left\vert A\right\rangle $, then it could duplicate \emph{any}
$\left\vert A\right\rangle $. Thus the bank can ensure security by creating
uniformly-random money states.

In what follows, let $\mathcal{S}$ be the set of all subspaces $A\leq
\mathbb{F}_{2}^{n}$ such that $\dim A=n/2$. Also, let $V_{A}^{\otimes
2}=\left(  \left\vert A\right\rangle \left\langle A\right\vert \right)
^{\otimes2}$ be the projector onto $\left\vert A\right\rangle ^{\otimes2}$.

\begin{theorem}
[Lower Bound for Average-Case Counterfeiting]\label{averagecase}Let
$A\leq\mathbb{F}_{2}^{n}$ be a uniformly-random element of $\mathcal{S}$.
\ Then given one copy of $\left\vert A\right\rangle $, as well as oracle
access to $U_{A^{\ast}}$, a counterfeiter $C$ needs $\Omega\left(
\sqrt{\varepsilon}2^{n/4}\right)  $ queries to prepare a $2n$-qubit state
$\rho$ that $V_{A}^{\otimes2}$ accepts with probability at least
$\varepsilon$, for all $1/\varepsilon=o\left(  2^{n/2}\right)  $. Here the
probability is taken over the choice of $A\in\mathcal{S}$, as well as the
behavior of $C$ and $V_{A}^{\otimes2}$.
\end{theorem}

\begin{proof}
Suppose we had a counterfeiter $C$ that violated the above. Using $C$ as a
black box, we will show how to construct a new counterfeiter $C^{\prime}%
$ that violates \expref{Corollary}{largeerrror}.

Given a (deterministically-chosen) money state $\left\vert A\right\rangle $
and oracle access to $U_{A^{\ast}}$, first choose an invertible linear map
$f:\mathbb{F}_{2}^{n}\rightarrow\mathbb{F}_{2}^{n}$ uniformly at random.
\ Then $f\left(  A\right)  $, the image of $A$ under $f$, is a
uniformly-random element of $\mathcal{S}$. Furthermore, the state
$\left\vert A\right\rangle $ can be transformed into $\left\vert f\left(
A\right)  \right\rangle $ straightforwardly, the oracle $U_{f\left(
A\right)  }$ can be simulated by composing $f$ with $U_{A}$, and the oracle
$U_{f\left(  A\right)  ^{\bot}}$ can likewise be simulated by composing
$f^{-T}$ with $U_{A}$ (where $f^{-T}$ denotes the inverse transpose of $f$).
\ So by using the counterfeiter $C$ for uniformly-random states, we can
produce a state $\rho_{f}$ that $V_{f\left(  A\right)  }^{\otimes2}$ accepts
with probability at least $\varepsilon$. By applying $f^{-1}$ to both
registers of $\rho_{f}$, we can then obtain a state $\rho$ that
$V_{A}^{\otimes2}$ accepts with probability at least $\varepsilon$, thereby
contradicting \expref{Corollary}{largeerrror}.
\end{proof}
\clearpage
\begingroup\raggedright
%
%
%
%
%
\providecommand{\bibhead}[1]{}
%
\expandafter\ifx\csname pdfbookmark\endcsname\relax%
  \providecommand{\tocrefpdfbookmark}{}
\else\providecommand{\tocrefpdfbookmark}{%
   \hypertarget{tocreferences}{}%
   \pdfbookmark[1]{References}{tocreferences}}%
\fi

\tocrefpdfbookmark
\begin{thebibliography}{10}

\bibitem{aar:col}\bibhead{aar:col}
{\sc Scott Aaronson}: Quantum lower bound for the collision problem.
\newblock In {\em Proc. 34th STOC}, pp. 635--642. ACM Press, 2002.
\newblock Preprint at \href{http://arxiv.org/abs/quant-ph/0111102}{arXiv}.
\newblock [\epfmtdoi{10.1145/509907.509999}]

\bibitem{aar:adv}\bibhead{aar:adv}
{\sc Scott Aaronson}: Limitations of quantum advice and one-way communication.
\newblock {\em Theory of Computing}, 1(1):1--28, 2005.
\newblock Preliminary version in
  \href{http://dx.doi.org/10.1109/CCC.2004.1313854}{CCC'04}. Preprint at
  \href{http://arxiv.org/abs/quant-ph/0402095}{arXiv}.
\newblock [\epfmtdoi{10.4086/toc.2005.v001a001}]

\bibitem{aar:qcopy}\bibhead{aar:qcopy}
{\sc Scott Aaronson}: Quantum copy-protection and quantum money.
\newblock In {\em Proc. 24th IEEE Conf. on Computational Complexity (CCC'09)},
  pp. 229--242. IEEE Comp. Soc. Press, 2009.
\newblock [\epfmtdoi{10.1109/CCC.2009.42}]

\bibitem{aar:private}\bibhead{aar:private}
{\sc Scott Aaronson}: On the security of private-key quantum money, 2013.
\newblock In preparation.

\bibitem{ak}\bibhead{ak}
{\sc Scott Aaronson and Greg Kuperberg}: Quantum versus classical proofs and
  advice.
\newblock {\em Theory of Computing}, 3(7):129--157, 2007.
\newblock Preliminary version in
  \href{http://dx.doi.org/10.1109/CCC.2007.27}{CCC'07}. Preprint at
  \href{http://arxiv.org/abs/quant-ph/0604056}{arXiv}.
\newblock [\epfmtdoi{10.4086/toc.2007.v003a007}]

\bibitem{aar:col-jacm}\bibhead{aar:col-jacm}
{\sc Scott Aaronson and Yaoyun Shi}: Quantum lower bounds for the collision and
  the element distinctness problems.
\newblock {\em J. ACM}, 51(4):595--605, 2004.
\newblock [\epfmtdoi{10.1145/1008731.1008735}]

\bibitem{ac:psn}\bibhead{ac:psn}
{\sc Martin~R. Albrecht and Carlos Cid}: Cold boot key recovery by solving
  polynomial systems with noise.
\newblock In {\em Proc. 9th Internat. Conf. on Applied Cryptography and Network
  Security (ACNS'11)}, pp. 57--72. Springer, 2011.
\newblock Available in \href{http://eprint.iacr.org/2011/038}{preprint}.
\newblock [\epfmtdoi{10.1007/978-3-642-21554-4\_4}]

\bibitem{alonclique}\bibhead{alonclique}
{\sc Noga Alon, Michael Krivelevich, and Benny Sudakov}: Finding a large hidden
  clique in a random graph.
\newblock {\em Random Structures {\&} Algorithms}, 13(3-4):457--466, 1998.
\newblock Preliminary version in
  \href{http://dl.acm.org/citation.cfm?id=314613.315014}{SODA'98}.
\newblock
  [\epfmtdoi{10.1002/(SICI)1098-2418(199810/12)13:3/4<457::AID-RSA14>3.0.CO;2-W}]

\bibitem{ambainis}\bibhead{ambainis}
{\sc Andris Ambainis}: Quantum lower bounds by quantum arguments.
\newblock {\em J. Comput. System Sci.}, 64(4):750--767, 2002.
\newblock Preliminary version in
  \href{http://dx.doi.org/10.1145/335305.335394}{STOC'00}. Preprint at
  \href{http://arxiv.org/abs/quant-ph/0002066}{arXiv}.
\newblock [\epfmtdoi{10.1006/jcss.2002.1826}]

\bibitem{amrr}\bibhead{amrr}
{\sc Andris Ambainis, Lo\"{i}ck Magnin, Martin Roetteler, and J\'{e}r\'{e}mie
  Roland}: Symmetry-assisted adversaries for quantum state generation.
\newblock In {\em Proc. 26th IEEE Conf. on Computational Complexity (CCC'11)},
  pp. 167--177. IEEE Comp. Soc. Press, 2011.
\newblock Preprint at \href{http://arxiv.org/abs/1012.2112}{arXiv}.
\newblock [\epfmtdoi{10.1109/CCC.2011.24}]

\bibitem{baraketal}\bibhead{baraketal}
{\sc Boaz Barak, Oded Goldreich, Russell Impagliazzo, Steven Rudich, Amit
  Sahai, Salil~P. Vadhan, and Ke~Yang}: On the (im)possibility of obfuscating
  programs.
\newblock {\em J. ACM}, 59(2):6, 2012.
\newblock Preliminary versions in
  \href{http://dx.doi.org/10.1007/3-540-44647-8_1}{CRYPTO'01} and
  \href{http://eccc.hpi-web.de/report/2001/057/}{ECCC}.
\newblock [\epfmtdoi{10.1145/2160158.2160159}]

\bibitem{bbcmw}\bibhead{bbcmw}
{\sc Robert Beals, Harry Buhrman, Richard Cleve, Michele Mosca, and Ronald
  de~Wolf}: Quantum lower bounds by polynomials.
\newblock {\em J. ACM}, 48(4):778--797, 2001.
\newblock Preliminary version in
  \href{http://dx.doi.org/10.1109/SFCS.1998.743485}{FOCS'98}. Preprint at
  \href{http://arxiv.org/abs/quant-ph/9802049}{arXiv}.
\newblock [\epfmtdoi{10.1145/502090.502097}]

\bibitem{bbbv}\bibhead{bbbv}
{\sc Charles~H. Bennett, Ethan Bernstein, Gilles Brassard, and Umesh~V.
  Vazirani}: Strengths and weaknesses of quantum computing.
\newblock {\em SIAM J. Comput.}, 26(5):1510--1523, 1997.
\newblock Preprint at \href{http://arxiv.org/abs/quant-ph/9701001}{arXiv}.
\newblock [\epfmtdoi{10.1137/S0097539796300933}]

\bibitem{bb84}\bibhead{bb84}
{\sc Charles~H. Bennett and Gilles Brassard}: Quantum cryptography: public key
  distribution and coin tossing.
\newblock In {\em Proc. IEEE Internat. Conf. on Computers, Systems and Signal
  Processing}, pp. 175--179, 1984.

\bibitem{bbbw}\bibhead{bbbw}
{\sc Charles~H. Bennett, Gilles Brassard, Seth Breidbart, and Stephen Wiesner}:
  Quantum cryptography, or unforgeable subway tokens.
\newblock In {\em Proceedings of CRYPTO}, pp. 267--275. Plenum Press, 1982.

\bibitem{bohr}\bibhead{bohr}
{\sc Niels Bohr}: {\em Atomic Physics and Human Knowledge}.
\newblock Dover, 2010.
\newblock First published 1961.

\bibitem{bdflsz}\bibhead{bdflsz}
{\sc Dan Boneh, \"{O}zg\"{u}r Dagdelen, Marc Fischlin, Anja Lehmann, Christian
  Schaffner, and Mark Zhandry}: Random oracles in a quantum world.
\newblock In {\em 17th Internat. Conf. on the Theory and Application of
  Cryptology and Information Security (ASIACRYPT'11)}, pp. 41--69, 2011.
\newblock Preprint at \href{http://arxiv.org/abs/1008.0931}{arXiv}.
\newblock [\epfmtdoi{10.1007/978-3-642-25385-0\_3}]

\bibitem{bffp:ip1s}\bibhead{bffp:ip1s}
{\sc Charles Bouillaguet, Jean-Charles Faug\`{e}re, Pierre-Alain Fouque, and
  Ludovic Perret}: Practical cryptanalysis of the identification scheme based
  on the isomorphism of polynomial with one secret problem.
\newblock In {\em 14th Internat. Conf. on Practice and Theory in Public Key
  Cryptography (PKC'11)}, pp. 473--493. Springer, 2011.
\newblock [\epfmtdoi{10.1007/978-3-642-19379-8\_29}]

\bibitem{bhmt}\bibhead{bhmt}
{\sc Gilles Brassard, Peter H{\o}yer, Michele Mosca, and Alain Tapp}: Quantum
  amplitude amplification and estimation.
\newblock In {\sc Jr. Samuel J.~Lomonaco and Howard~E. Brandt}, editors, {\em
  Quantum Computation and Information}, Contemporary Mathematics Series. AMS,
  2002.
\newblock Preprint at \href{http://arxiv.org/abs/quant-ph/0005055}{arXiv}.

\bibitem{cash}\bibhead{cash}
{\sc David Cash, Dennis Hofheinz, Eike Kiltz, and Chris Peikert}: Bonsai trees,
  or how to delegate a lattice basis.
\newblock {\em J. Cryptology}, 25(4):601--639, 2012.
\newblock Preliminary version in
  \href{http://dx.doi.org/10.1007/978-3-642-13190-5_27}{EUROCRYPT'10}.
\newblock [\epfmtdoi{10.1007/s00145-011-9105-2}]

\bibitem{crr}\bibhead{crr}
{\sc Sourav Chakraborty, Jaikumar Radhakrishnan, and Nandakumar Raghunathan}:
  Bounds for error reduction with few quantum queries.
\newblock In {\em Proc. 9th Internat. Workshop on Randomization and Computation
  (RANDOM'05)}, pp. 245--256. Springer, 2005.
\newblock [\epfmtdoi{10.1007/11538462\_21}]

\bibitem{dy:mpkc}\bibhead{dy:mpkc}
{\sc Jintai Ding and Bo-Yin Yang}: Multivariate public key cryptography.
\newblock In {\sc Daniel~J. Bernstein, Johannes Buchmann, and Erik Dahmen},
  editors, {\em Post-Quantum Cryptography}, pp. 193--241. Springer, 2009.
\newblock [\epfmtdoi{10.1007/978-3-540-88702-7\_6}]

\bibitem{farhi:restore}\bibhead{farhi:restore}
{\sc Edward Farhi, David Gosset, Avinatan Hassidim, Andrew Lutomirski, Daniel
  Nagaj, and Peter Shor}: Quantum state restoration and single-copy tomography
  for ground states of {H}amiltonians.
\newblock {\em Phys. Rev. Lett.}, 105(19):190503, 2010.
\newblock Preprint at \href{http://arxiv.org/abs/0912.3823}{arXiv}.
\newblock [\epfmtdoi{10.1103/PhysRevLett.105.190503}]

\bibitem{knots}\bibhead{knots}
{\sc Edward Farhi, David Gosset, Avinatan Hassidim, Andrew Lutomirski, and
  Peter~W. Shor}: Quantum money from knots.
\newblock In {\em Proc. 3rd Innovations in Theoretical Comput. Sci. (ITCS'12)},
  pp. 276--289. ACM Press, 2012.
\newblock Preprint at \href{http://arxiv.org/abs/1004.5127}{arXiv}.
\newblock [\epfmtdoi{10.1145/2090236.2090260}]

\bibitem{gavinsky:money}\bibhead{gavinsky:money}
{\sc Dmitry Gavinsky}: Quantum money with classical verification.
\newblock In {\em Proc. 27th IEEE Conf. on Computational Complexity (CCC'12)},
  pp. 42--52. IEEE Comp. Soc. Press, 2012.
\newblock Preprint at \href{http://arxiv.org/abs/1109.0372}{arXiv}.
\newblock [\epfmtdoi{10.1109/CCC.2012.10}]

\bibitem{gms:ip1s}\bibhead{gms:ip1s}
{\sc Willi Geiselmann, Willi Meier, and Rainer Steinwandt}: An attack on the
  isomorphisms of polynomials problem with one secret.
\newblock {\em Int. J. Inf. Sec.}, 2(1):59--64, 2003.
\newblock [\epfmtdoi{10.1007/s10207-003-0025-5}]

\bibitem{grover}\bibhead{grover}
{\sc Lov~K. Grover}: A fast quantum mechanical algorithm for database search.
\newblock In {\em Proc. 28th STOC}, pp. 212--219. ACM Press, 1996.
\newblock Preprint at \href{http://arxiv.org/abs/quant-ph/9605043}{arXiv}.
\newblock [\epfmtdoi{10.1145/237814.237866}]

\bibitem{hl:tnp}\bibhead{hl:tnp}
{\sc Christopher~J. Hillar and Lek-Heng Lim}: Most tensor problems are
  {NP}-hard.
\newblock Technical report, 2009.
\newblock [\epfmt{arxiv}{0911.1393}]

\bibitem{lmrss}\bibhead{lmrss}
{\sc Troy Lee, Rajat Mittal, Ben~W. Reichardt, Robert Spalek, and Mario
  Szegedy}: Quantum query complexity of state conversion.
\newblock In {\em Proc. 52nd FOCS}, pp. 344--353. IEEE Comp. Soc. Press, 2011.
\newblock Preprint at \href{http://arxiv.org/abs/1011.3020}{arXiv}.
\newblock [\epfmtdoi{10.1109/FOCS.2011.75}]

\bibitem{lutomirski:attack}\bibhead{lutomirski:attack}
{\sc Andrew Lutomirski}: An online attack against {W}iesner's quantum money.
\newblock Technical report, 2010.
\newblock [\epfmt{arxiv}{1010.0256}]

\bibitem{lutomirski:scp}\bibhead{lutomirski:scp}
{\sc Andrew Lutomirski}: Component mixers and a hardness result for
  counterfeiting quantum money.
\newblock Technical report, 2011.
\newblock [\epfmt{arxiv}{1107.0321}]

\bibitem{breaking}\bibhead{breaking}
{\sc Andrew Lutomirski, Scott Aaronson, Edward Farhi, David Gosset, Jonathan~A.
  Kelner, Avinatan Hassidim, and Peter~W. Shor}: Breaking and making quantum
  money: Toward a new quantum cryptographic protocol.
\newblock In {\em Proc. Innovations in Comput. Sci. (ICS'10)}, pp. 20--31.
  Tsinghua University Press, 2010.
\newblock
  \href{http://conference.itcs.tsinghua.edu.cn/ICS2010/papers.html}{ICS'10}.
  Preprint at \href{http://arxiv.org/abs/0912.3825}{arXiv}.

\bibitem{molina}\bibhead{molina}
{\sc Abel Molina, Thomas Vidick, and John Watrous}: Optimal counterfeiting
  attacks and generalizations for {W}iesner's quantum money.
\newblock In {\em Proc. 7th Conf. Theory of Quantum Computation, Communication,
  and Cryptography (TQC'12)}, pp. 45--64, 2012.
\newblock Preprint at \href{http://arxiv.org/abs/1202.4010}{arXiv}.
\newblock [\epfmtdoi{10.1007/978-3-642-35656-8\_4}]

\bibitem{moscastebila}\bibhead{moscastebila}
{\sc Michele Mosca and Douglas Stebila}: Quantum coins.
\newblock In {\em Error-Correcting Codes, Finite Geometries and Cryptography},
  volume 523, pp. 35--47. Amer. Math. Soc., 2010.
\newblock Preprint at \href{http://arxiv.org/abs/0911.1295}{arXiv}.

\bibitem{naoryung}\bibhead{naoryung}
{\sc Moni Naor and Moti Yung}: Universal one-way hash functions and their
  cryptographic applications.
\newblock In {\em Proc. 21st STOC}, pp. 33--43. ACM Press, 1989.
\newblock [\epfmtdoi{10.1145/73007.73011}]

\bibitem{nc}\bibhead{nc}
{\sc Michael Nielsen and Isaac Chuang}: {\em Quantum Computation and Quantum
  Information}.
\newblock Cambridge University Press, 2000.
\newblock [\epfmt{acm}{544199}]

\bibitem{pastawski}\bibhead{pastawski}
{\sc Fernando Pastawski, Norman~Y. Yao, Liang Jiang, Mikhail~D. Lukin, and
  J.~Ignacio Cirac}: Unforgeable noise-tolerant quantum tokens.
\newblock Technical report, 2011.
\newblock [\epfmt{arxiv}{1112.5456}]

\bibitem{cgp:ip1s}\bibhead{cgp:ip1s}
{\sc Jacques Patarin, Louis Goubin, and Nicolas Courtois}: Improved algorithms
  for isomorphisms of polynomials.
\newblock In {\em Proc. Internat. Conf. on the Theory and Application of
  Cryptographic Techniques (EUROCRYPT'98)}, pp. 184--200. Springer, 1998.
\newblock [\epfmtdoi{10.1007/BFb0054126}]

\bibitem{regev}\bibhead{regev}
{\sc Oded Regev}: On lattices, learning with errors, random linear codes, and
  cryptography.
\newblock {\em J. ACM}, 56(6):34, 2009.
\newblock Preliminary version in
  \href{http://dx.doi.org/10.1145/1060590.1060603}{STOC'05}.
\newblock [\epfmtdoi{10.1145/1568318.1568324}]

\bibitem{rompel}\bibhead{rompel}
{\sc John Rompel}: One-way functions are necessary and sufficient for secure
  signatures.
\newblock In {\em Proc. 22nd STOC}, pp. 387--394. ACM Press, 1990.
\newblock [\epfmtdoi{10.1145/100216.100269}]

\bibitem{tgp}\bibhead{tgp}
{\sc Tathagat Tulsi, Lov~K. Grover, and Apoorva Patel}: A new algorithm for
  fixed point quantum search.
\newblock {\em Quantum Inf. Comput.}, 6(6):483--494, 2006.
\newblock \href{http://www.rintonpress.com/journals/qiconline.html\#v6n6}{QIC}.
\newblock [\epfmt{acm}{2011693}]

\bibitem{wiesner}\bibhead{wiesner}
{\sc Stephen Wiesner}: Conjugate coding.
\newblock {\em SIGACT News}, 15(1):78--88, 1983.
\newblock Original manuscript written circa 1970.
\newblock [\epfmtdoi{10.1145/1008908.1008920}]

\end{thebibliography}
\endgroup
\end{document}
